\documentclass[../../../include/open-logic-section]{subfiles}

\begin{document}
	
\olfileid{sth}{z}{nat}
\olsection{Milih Wilangan Asli Kita}

%Upamane kita nampa yen prinsip informal \stagesinf{}
%bisa diwenehi dhasar. Nanging aksioma tartamtu kang dijupuk
%saka prinsip informal iku uga perlu dirembug.

Ing \olref[infinity-again]{defnomega}, kita wis netepake
kanthi cetha teges ungkapan ``wilangan asli''. Kepiye
panjenengan kudu mangerteni panetepan iki? Iki dudu
pratelan metafisik, nanging mung kaputusan kanggo
\emph{nganggep} sawenehing himpunan minangka wilangan asli.
Kita wis nyinggung alasan-alasane ing
\olref[sfr][rel][ref]{sec}, \olref[sfr][arith][ref]{sec}, lan
\olref[sfr][infinite][dedekindsproof]{sec}. Nanging alasan
iku bisa kita andharake kanthi luwih cetha maneh.

Aksioma Tanpa Wates kang kita gunakake manut
\citet{VonNeumann1925}. Nanging ana aksioma liyane kang
bisa kita pilih minangka gantine:

\begin{defish}
\emph{Aksioma Tanpa Wates Zermelo saka
\citeyear{Zermelo1908Untersuchungen}.} Ana himpunan $A$
kanthi $\emptyset \in A$ lan $(\forall x \in A)\{x\} \in A$.
\end{defish}

Yen kita nggunakake aksioma Zermelo minangka ganti
Aksioma Tanpa Wates kita (kang diilhami von Neumann),
kita uga bakal oleh himpunan tanpa wates miturut Dedekind,
mula uga aljabar Dedekind. Miturut cara Zermelo, unsur
kang dipilih ing aljabar kita tetep $\emptyset$ (wakil
kita kanggo $0$), nanging injeksine diwenehake dening
pemetaan $x \mapsto \{x\}$, dudu $x \mapsto x \cup \{x\}$.
Akibat paling gampang yaiku Zermelo nganggep $2$ minangka
$\{\{\emptyset\}\}$, dene kita (manut von Neumann)
nganggep $2$ minangka $\{\emptyset, \{\emptyset\}\}$.

Apa sebabe kita milih salah siji Aksioma Tanpa Wates?
Alasan praktis kang utama yaiku cara von Neumann bisa
ditrapake kanthi becik uga marang wilangan transfinit.
Kita bakal nliti iki wiwit
\olref[ordinals][]{chap}. Nanging saka sudut pandang
prasaja kanggo \emph{nindakake aritmetika}, loro-lorone
padha migunani. Mula yen ana wong kandha yen wilangan
asli \emph{iku} himpunan, pitakonan kang lumrah yaiku:
\emph{Himpunan endi kang dadi wilangan asli?}

Pitakonan kang trep iki dadi misuwur amarga
\citet{Benacerraf1965}. Nanging kudu dielingi yen iki
namung conto paling misuwur saka kedadeyan kang wis bola-bali
kita temoni. Pokoke mangkene: teori himpunan menehi cara
kanggo \emph{nyimulasi} maneka jinis entitas ``intuitif'':
wilangan real, rasional, bulat, lan asli; uga pasangan
urut, fungsi, lan relasi. Nanging teori himpunan ora tau
menehi pilihan simulasi kang \emph{tunggal}. \emph{Mesthi}
ana pilihan liyane kang kanthi gampang bisa nyukupi
kabutuhan kita kanthi padha becike.

\end{document}
